% Lokalisasi khusus permukaan pembaca untuk formulir ujian dalam edisi lengkap.
% Nama perintah sumber dipertahankan; hanya teks yang dicetak yang dilokalkan.

\renewcommand{\fachbereich}{Bidang studi}
\renewcommand{\dozent}{Pengajar}
\renewcommand{\klausurdatum}{Tanggal}
\renewcommand{\klausurgebiet}{Matematika}
\renewcommand{\klausurtyp}{Ujian}

\renewcommand{\klausurmodus}{%
Durasi: sepuluh menit untuk membaca dan dua jam penuh untuk mengerjakan.
Alat bantu tidak diperbolehkan. Semua jawaban harus diberi alasan.

Jumlah poin seluruhnya $64$. Berlaku aturan ambang dasar, yakni penilaian
untuk setiap soal atau bagian soal pada umumnya dimulai pada separuh jumlah
poin. Untuk lulus diperlukan $16$ poin; mulai $32$ poin diperoleh nilai satu.

Tuliskan nama dan nomor mahasiswa Anda dengan jelas pada lembar sampul dan
setiap lembar berikutnya. Semoga berhasil!}

\renewcommand{\testklausurmodus}{\klausurmodus}

\renewcommand{\klausurname}{%
\renewcommand{\arraystretch}{1.5}%
\begin{tabular}{@{}p{3.2cm}p{10cm}}
Nama lengkap: &
..................................................................................\\
Nomor mahasiswa: &
..................................................................................\\
\end{tabular}}

\renewcommand{\klausurvorspann}[5]{%
\begin{tabular}{@{}p{3.2cm}p{10cm}}
Bidang studi: & #1\\
Tanggal: & #2\\
Pengajar: & #3
\end{tabular}
\vspace{3ex}

\begin{center}
{\bfseries\large Ujian\\#4\ #5}
\end{center}

\vspace{2ex}
\klausurmodus
\vspace{2ex}
\klausurname\par}

\renewcommand{\testklausurvorspann}[5]{\klausurvorspann{#1}{#2}{#3}{#4}{#5}}
\renewcommand{\klausurmitloesungenvorspann}[5]{%
\klausurvorspann{#1}{#2}{#3}{#4}{#5}
\begin{center}\bfseries Formulir dengan solusi\end{center}}

\renewcommand{\klausurnote}{%
\begin{tabular}{@{}p{3.2cm}p{10cm}}Nilai: & \end{tabular}}

\renewcommand{\klausurtaeuschungwarnung}{%
\vspace{2ex}%
Saya memahami bahwa pelanggaran berat terhadap aturan ujian dapat membatalkan
seluruh komponen penilaian dalam mata kuliah ini. Penggunaan literatur atau
alat bantu elektronik termasuk pelanggaran berat.
\vspace{1ex}}

\renewcommand{\klausureinverstaendnis}{%
\vspace{2ex}
Dengan tanda tangan saya, saya menyetujui bahwa hasil ujian dapat diumumkan
bersama nomor mahasiswa saya.
\vspace{1ex}

\begin{tabular}{@{}p{3.2cm}p{10cm}}
Tanda tangan: &
.....................................................................................
\end{tabular}}

\renewcommand{\inputaufgabeklausurloesung}[3]{%
\par\newpage
\begin{aufgabe}
{\ifthenelse{\equal{#1}{}}{}{\ifthenelse{\equal{#1}{1}}{(1 poin)}{(#1 poin)}}}%
\par\bigskip #2
\end{aufgabe}
\underline{Solusi:}\par\bigskip #3}

\renewcommand{\inputaufgabeloesung}[2]{%
\begin{aufgabe}#1\end{aufgabe}\aufgabenachskip
\bigskip
Solusi
\bigskip #2}

\renewcommand{\inputaufgabeloesungvar}[2]{%
\aufgabevorskip Soal: #1\aufgabenachskip Solusi: #2}

\renewcommand{\inputaufgabepunkteloesung}[3]{%
\aufgabevorskip
\begin{aufgabe}
{\ifthenelse{\equal{#1}{}}{}{\ifthenelse{\equal{#1}{1}}{(1 poin)}{(#1 poin)}}}%
\aufgabepunktskip #2
\end{aufgabe}
\aufgabenachskip\loesung{#3}}

\renewcommand{\inputaufgabepunktenummervonhand}[3]{%
\aufgabevorskip{\bfseries Soal #1} (#2 poin)\aufgabepunktskip #3\aufgabenachskip}

\renewcommand{\loesung}[1]{\underline{Solusi:} #1}

% Delapan formulir sumber memakai huruf penampung K sebagai nilai penghitung
% bagian. Pertahankan byte formulir dan tafsirkan penampung itu sebagai awal
% penomoran hanya pada lapisan kompilasi pembaca.
\let\OExamOriginalSetCounter\setcounter
\renewcommand{\setcounter}[2]{%
\ifthenelse{\equal{#1}{bagian}}%
  {\OExamOriginalSetCounter{section}{#2}}%
  {\ifthenelse{\equal{#1}{section}}%
  {\ifthenelse{\equal{#2}{K}}{\OExamOriginalSetCounter{section}{0}}{\OExamOriginalSetCounter{#1}{#2}}}%
  {\OExamOriginalSetCounter{#1}{#2}}}}

\newcommand{\OExamPointTable}[4]{%
\vspace{3ex}
\begin{center}
\begingroup
\scriptsize
\renewcommand{\arraystretch}{1.45}%
\setlength{\tabcolsep}{0.35ex}%
\begin{tabular}{@{}|l|*{#1}{c|}c|}
\hline
Soal & #2 & Jumlah\\ \hline
Maks. & #3 & #4\\ \hline
Diperoleh & \multicolumn{#1}{c|}{} & \\ \hline
\end{tabular}
\endgroup
\end{center}
\vspace{2ex}}

\renewcommand{\punktetabelleelf}{%
\OExamPointTable{11}%
{1&2&3&4&5&6&7&8&9&10&11}%
{\aeins&\azwei&\adrei&\avier&\afuenf&\asechs&\asieben&\aacht&\aneun&\azehn&\aelf}%
{\azwoelf}}

\renewcommand{\punktetabelledreizehn}{%
\OExamPointTable{13}%
{1&2&3&4&5&6&7&8&9&10&11&12&13}%
{\aeins&\azwei&\adrei&\avier&\afuenf&\asechs&\asieben&\aacht&\aneun&\azehn&\aelf&\azwoelf&\adreizehn}%
{\avierzehn}}

\renewcommand{\punktetabellevierzehn}{%
\OExamPointTable{14}%
{1&2&3&4&5&6&7&8&9&10&11&12&13&14}%
{\aeins&\azwei&\adrei&\avier&\afuenf&\asechs&\asieben&\aacht&\aneun&\azehn&\aelf&\azwoelf&\adreizehn&\avierzehn}%
{\afuenfzehn}}

\renewcommand{\punktetabellesechzehn}{%
\OExamPointTable{16}%
{1&2&3&4&5&6&7&8&9&10&11&12&13&14&15&16}%
{\aeins&\azwei&\adrei&\avier&\afuenf&\asechs&\asieben&\aacht&\aneun&\azehn&\aelf&\azwoelf&\adreizehn&\avierzehn&\afuenfzehn&\asechzehn}%
{\asiebzehn}}

\renewcommand{\punktetabellesiebzehn}{%
\OExamPointTable{17}%
{1&2&3&4&5&6&7&8&9&10&11&12&13&14&15&16&17}%
{\aeins&\azwei&\adrei&\avier&\afuenf&\asechs&\asieben&\aacht&\aneun&\azehn&\aelf&\azwoelf&\adreizehn&\avierzehn&\afuenfzehn&\asechzehn&\asiebzehn}%
{\aachtzehn}}
